% !TEX root = ../main-es.tex
%TC:macro \limite [ignore,ignore]
%TC:macro \nota [ignore]
%TC:macro \TODO [ignore]
%TC:macro \parencite [ignore]
%TC:macro \textcite [ignore]
\section{Discusión}\label{sec:discusion}

\subsection{Por qué la frontera tiene esa forma}

La sesión del curso sobre el modelo predador–presa explica la estabilidad de ese sistema diagnosticando primero por qué colapsa y enumerando después los mecanismos que lo sostienen, y la misma estructura sirve para leer la frontera $(\lambda, c)$ como un sistema y no figura por figura, con la advertencia de que el paralelo es una \emph{analogía} y no un mecanismo, pues el modelo no tiene poblaciones acopladas ni ciclos (la dinámica no mostró ninguna oscilación). Cada pieza se rotula entonces como mecanismo o como analogía, y la evidencia sale de los barridos.

\begin{enumerate}[leftmargin=1.6em, itemsep=0.3em]
  \item \textbf{El diagnóstico, deliberar se come el tiempo} (mecanismo). Con $c=0$ el lazo abierto es la mejor política o empata en todo $\lambda$, pero desde $c=2$ y $\lambda\ge0.25$ la mejor pasa a ser \MB, porque en la celda $\lambda=0.5$, $c=5$ el lazo cerrado pasa 282.2 de sus 400 ticks deliberando y el abierto 153.6, y ese impuesto se cobra desde el primer tramo de la corrida. El frente que importa con dinamismo alto deja de ser «comprometerse frente a reparar» y pasa a ser «planificar frente a no planificar».
  \item \textbf{La observabilidad, un desplazador y no el diagnóstico} (mecanismo). El radio mueve sobre todo el nivel de todas las políticas, y la creencia obsoleta existe (el oráculo no puede acertar sin ver el objetivo), pero no es lo que hunde a las políticas con plan, pues las preguntas lógicas que la atacan (descartar candidatos vistos vacíos y elegir objetivo por el peor caso) no compran nada medible por sí solas, según la ablación con \ahora, y lo único que ayuda es dejar de pagar reparaciones, que es otra vez el punto 1.
  \item \textbf{La respuesta con retardo} (analogía con la reproducción lenta del predador; el retardo en sí es mecanismo). Deliberar ocupa $c$ ticks y el plan se calcula al empezar, de modo que llega desactualizado si el mundo se movió mientras tanto. La analogía se rompe justamente aquí, pues en Lotka–Volterra el retardo produce ciclos, mientras que en este modelo no hay retroalimentación del agente sobre el mundo que pueda producirlos, y el retardo se traduce en un impuesto de ritmo constante: Las entregas acumuladas y los ticks deliberando crecen en línea recta.
  \item \textbf{El \emph{thrashing} de replanificación} (mecanismo). Cada divergencia nueva interrumpe la deliberación en curso y los ticks gastados se pierden, lo cual con $r=\infty$, $c=5$ y $\lambda=0.5$ deja al lazo cerrado en utilidad 0.0, frente a 34.0 del abierto, que también delibera tras cada fallo pero no se deja interrumpir por las divergencias.
  \item \textbf{Los amortiguadores} (analogía con el tope de energía del predador). Un presupuesto pequeño recorta el \emph{thrashing} y el horizonte lógico también, pero con dinamismo alto ninguno supera a no reparar de forma proactiva ($B=0$ rinde 27.0 frente a 23.7 con $B=3$, y \KB{} rinde lo mismo que el lazo abierto, \dif{+4.3}{5.3}), de forma que el mejor amortiguador medido es no reparar al ver la divergencia.
  \item \textbf{La capacidad de carga} (analogía con el pasto que se repone). Cada entrega repone un recurso, de forma que el mundo nunca se queda sin trabajo, y la abundancia cambia el valor de replanificar, aunque al revés de lo esperado. La analogía vale solo en que la abundancia fija la escala, pues el mecanismo es el costo de oportunidad (§\ref{sec:disc-oportunidad}).
\end{enumerate}

Dicho en una frase coherente con los datos: Cuando reparar cuesta, el lazo abierto y el cerrado quedan cerca y, con el entorno cambiando rápido, gana el agente que no delibera, tanto más cuanto más abundan los objetivos; la observabilidad parcial mueve sobre todo el nivel; el \emph{thrashing} es real y un presupuesto pequeño o un horizonte lógico lo recortan, aunque con dinamismo alto no reparar de forma proactiva lo recorta igual; y reparar en caliente frente a reiniciar en frío compensa en casi todo el rango con $\lambda>0$, salvo donde ambas políticas colapsan.

\subsection{Un criterio operable de cuándo reescribir}\label{sec:disc-criterio}

La pregunta del trabajo pedía condiciones del entorno, y la respuesta que dan los datos se puede convertir en un procedimiento. El parámetro que hay que medir es $\lambda$, porque es el factor que más mueve a toda política con plan (hasta el 62\,\% de la varianza) y el que decide entre planificar y no planificar, mientras que la competencia $p$ del replanificador pesa poco ($\eta^2\le0.015$). Además, $\lambda$ se puede medir, pues la calibración por ABC lo recupera a partir de la telemetría de un motor que corre en frío (el valor verdadero cayó dentro del intervalo creíble del 90\,\% en los tres casos), aunque con precisión moderada, mientras que $p$ no es identificable con esa telemetría, ya que con un motor que solo replanifica al fallar los errores del replanificador se confunden con el dinamismo, y medirla exigiría etiquetar cada reparación contra la óptima, cosa que esa telemetría no registra. Por suerte, la decisión que más importa depende de $\lambda$ y no de $p$.

Con $\lambda$ estimado y el costo $c$ de cada reparación conocido (en un motor real, la latencia durante la cual el mundo sigue cambiando), la recomendación se lee en la frontera en dos etapas, primero si planificar y después, si conviene, cómo. El cuadro~\ref{tab:criterio} la resume para el régimen barrido ($r=3$, $p=0.9$), con la advertencia de que las unidades de $\lambda$ y $c$ son las del modelo (probabilidad de deriva por tick y ticks por reparación), de forma que llevarlas a un sistema real exige fijar qué es un tick en él.

\begin{table}[htbp]
\centering
\footnotesize
\begin{tabular}{@{}L{3.4cm} L{4.4cm} L{7.0cm}@{}}
\toprule
\textbf{Régimen} & \textbf{Recomendación} & \textbf{Evidencia} \\
\midrule
$\lambda\approx0$ & Planificar & Las políticas con plan casi duplican a las reactivas (93.3 contra 57.0) \\ \sepfilas
$\lambda$ entre 0.05 y 0.1, $c\le5$ & Planificar y reparar en caliente; monitorear o no es indiferente & \MB{} pierde por entre 17 y 44 puntos; caliente supera a frío desde $\lambda=0.1$; \OL$-$\CL{} dentro del ruido \\ \sepfilas
$\lambda\approx0.25$, $c=0$ & Planificar, en lazo abierto & \OL{} 111.7 contra 86.0 del cerrado y de \MB \\ \sepfilas
$\lambda\approx0.25$, $c=2$; $\lambda\approx0.1$, $c=10$ & Zona de indiferencia & \MB{} frente al mejor planificador: \dif{+4.7}{9.4} y \dif{+4.7}{8.8} \\ \sepfilas
$\lambda\approx0.25$, $c\ge5$ & No planificar & \MB{} gana por \dif{+30.0}{10.1} y \dif{+52.3}{9.2} \\ \sepfilas
$\lambda\approx0.5$, $c=0$ & Planificar o no, casi indiferente & \OL{} 120.3 contra 98.3 de \MB{} (\dif{-22.0}{11.2}) \\ \sepfilas
$\lambda\approx0.5$, $c\ge2$ & No planificar & \MB{} gana por entre $+45$ y $+86$ \\ \sepfilas
Visión total, $c=5$, $\lambda\ge0.25$ & Si se planifica, no monitorear & \OL{} supera a \CL{} por \dif{+20.0}{5.1} y \dif{+34.0}{4.5} \\
\bottomrule
\end{tabular}
\caption{Lectura de la frontera como criterio de decisión, con $r=3$ y $p=0.9$ salvo la última fila. Con más objetivos simultáneos la región de «no planificar» se extiende hacia $\lambda$ más bajos (§\ref{sec:res-escasez}).}
\label{tab:criterio}
\end{table}

Llevado a las tres decisiones de diseño que motivaron el trabajo, el criterio dice lo siguiente. Primeramente, construir la reparación en caliente compensa en cuanto el entorno cambia con alguna frecuencia, pues el reinicio en frío pierde entre 16 y 22 puntos desde $\lambda=0.1$, salvo que deliberar sea tan caro que el agente deje de completar tareas, mientras que en un entorno casi estático reiniciar en frío no cuesta nada medible. Progresivamente, un tope fijo de iteraciones no tiene respaldo en los datos como la razón por la que un agente funciona, y lo que corresponde es derivar cuánto deliberar del dinamismo estimado y del costo de cada reparación, lo cual con dinamismo alto y deliberación cara equivale a no deliberar. Finalmente, un presupuesto de ediciones, del que la autorización humana es una versión manual, evita el peor caso cuando el entorno cambia rápido, pero no rescata al agente, de modo que como freno contra el \emph{thrashing} rinde menos que la decisión, anterior, de no reparar al ver cada divergencia.

\subsection{Lógica frente a probabilidad}

Los dos agentes que infieren dejan un resultado con dos niveles. En el de la \emph{creencia}, el puente de la lógica a la probabilidad funciona sin pérdida: El soporte del filtro coincide tick a tick con la región lógica (lo que la KB considera desconocido es exactamente lo que tiene probabilidad positiva), $\lambda$ se aprende mirando y el filtro está calibrado. En el de la \emph{decisión}, en cambio, tener la probabilidad correcta no bastó, porque una regla de utilidad esperada es tan buena como el modelo de costos con que se evalúa, y la regla $c<(1-q)(c+k)$ omite que el plan reparado nace desactualizado y puede interrumpirse, de forma que repara de más con dinamismo alto. La lógica gana entonces por ser más tacaña, ya que solo repara cuando el recurso ya no puede estar en la celda del plan, y esa tacañería es justamente lo que el experimento premia, pues todo lo que gana \KB{} sobre el lazo cerrado viene de no pagar reparaciones (la ablación con \ahora{} lo deja en cero). Dicho de otra forma, el experimento no refuta que la probabilidad sirva para decidir, sino que muestra que el cuello de botella se movió de \emph{qué sé} a \emph{cuánto cuesta lo que hago con lo que sé}. Una regla corregida, que cargara a la reparación la probabilidad de que el plan reparado caduque durante los $c$ ticks que tarda, es la continuación natural, y no se probó.

El agente lógico deja además una respuesta a la pregunta que la sesión de agentes lógicos planteaba sobre qué pasa con la monotonía si el mundo se mueve: La monotonía no se rompe, se indexa por el tiempo. Ninguna oración se retira nunca, pues $\mathit{Visto}(r,c,t)$ sigue siendo verdadera; lo que caduca es la conclusión «está ahí ahora», que pasa de afirmarse a derivarse con el axioma de deriva, y la prueba de solidez dice que esa derivación nunca deja fuera la celda real. La tabla de creencias de las demás políticas, en cambio, borra y sobrescribe sin dejar rastro, y por eso no puede saber cuán vieja es una creencia. Vale la pena anotar, finalmente, la advertencia de H3 para los agentes basados en modelos de lenguaje, en los que «preguntarse si el plan sigue siendo válido» suele ser otra llamada al modelo: En este experimento, razonar solo compensa si es casi gratis frente a deliberar, pues un solo tick por consulta basta para invertir la ventaja.

\subsection{El costo de oportunidad}\label{sec:disc-oportunidad}

La hipótesis de escasez se escribió con una intuición razonable, que con objetivos escasos cada uno vale más y protegerlo con reparaciones debería compensar, y los datos la refutaron en los dos regímenes. La explicación que la propia hipótesis dejaba escrita para el caso contrario es la que se sostiene: Con objetivos escasos el cuello de botella es encontrar y alcanzar un recurso, y un tick deliberando cuesta poco porque de todos modos no había otra entrega cerca, mientras que con objetivos abundantes cada tick quieto es una entrega que se deja de hacer. El mecanismo es doble, pues con más objetivos hay más divergencias posibles y más reparaciones que pagar, justo cuando cada tick vale más. Donde deliberar es gratis no hay costo de oportunidad, y por eso en ese régimen la abundancia favorece a quien repara. La lectura para el diseño es que el valor de la deliberación no es una propiedad del agente ni del entorno por separado, sino de la relación entre lo que cuesta pensar y lo que se deja de hacer mientras se piensa.

\subsection{Relación con la agenda de los agentes auto-evolutivos}

La dimensión \emph{when} de la agenda de los agentes auto-evolutivos \parencite{gao2026selfevolving} pregunta cuándo debe un agente modificarse, y este trabajo ofrece una respuesta acotada y medible para el caso de la reescritura del plan: No conviene reescribir cada vez que el plan y el mundo dejan de coincidir, sino según cuánto cambia el entorno y cuánto cuesta cada reescritura, y con dinamismo alto y reescritura cara conviene no planificar en absoluto. Esa respuesta matiza, en lugar de contradecir, la presentación de \textcite[§11.5]{russellNorvig2020} del monitoreo de ejecución y la replanificación como la respuesta a un mundo que no se comporta como el plan supone, pues monitorear compensa cuando reparar es barato y bueno (el control 2) y deja de compensar en cuanto reparar cuesta tiempo en un mundo que no espera. Es también coherente con que la línea formal sobre procesos parcialmente observables se pregunte cuándo es seguro saltarse una replanificación \parencite{kong2026open}, y le da la razón a \textcite{brooks1991intelligence} en una región del espacio de parámetros, con dos salvedades: El ganador no es puramente reactivo, pues mantiene creencias, y lo que aparece no es una postura general sino una región delimitada por $\lambda$ y $c$.

Para los sistemas de la clase, la abstracción del oráculo deja una lección que sobrevive al cambio de modelo de lenguaje, pues de los dos parámetros del replanificador el que decide es el costo y no la competencia: Una reescritura casi perfecta que tarda lo suficiente para que el mundo cambie durante ella rinde menos que no reescribir, mientras que mejorar la competencia de un replanificador ya razonable mueve poco el resultado. Cabe mencionar que todo esto se midió con $\kappa=0$, es decir, contando solo el tiempo perdido, de forma que con un costo explícito por tokens la frontera se desplazaría todavía más en contra de deliberar, en una magnitud que este trabajo no midió.
